% ---------------------------------------------------------------------------
% Introducción (norma 3.12)
% ---------------------------------------------------------------------------
% NO se numera como capítulo (norma 2.6). El primer capítulo numerado es el que
% sigue inmediatamente a esta sección.
%
% La norma exige que contenga, en este orden: antecedentes; justificación e
% importancia del trabajo, con la hipótesis de manera implícita; planteamiento
% del problema; objetivo general; y objetivos específicos.
%
% Redactada en tercera persona. Los acrónimos se explican en su primera
% aparición dentro de esta sección, con su traducción cuando proceden de otro
% idioma (norma 2.4).

\seccionSinNumerar{Introducción}

El desarrollo de firmware para microcontroladores de 8 bits exige un
conocimiento detallado de la hoja de datos del dispositivo, la realización de
proyectos de prueba que permitan comprender la configuración de cada periférico
y el dominio de un entorno de desarrollo específico para cada familia empleada.
Ese conocimiento no es transferible entre fabricantes ni, en muchos casos, entre
familias de un mismo fabricante, de modo que cada nuevo dispositivo impone un
período de aprendizaje antes de que resulte productivo.

El presente trabajo, realizado como pasantía larga en la empresa Domotai C.A.,
aborda ese coste mediante un módulo de generación de código asistida por
inteligencia artificial para el microcontrolador PIC18F47Q10 de Microchip
Technology.

\subsection*{Antecedentes}

Los modelos de lenguaje de gran escala, designados habitualmente por sus siglas
en inglés LLM (\emph{Large Language Model}), han alcanzado en los últimos años
una capacidad notable para producir código a partir de descripciones en lenguaje
natural. El trabajo que estableció el marco de evaluación hoy dominante presentó
un modelo entrenado sobre código junto con un conjunto de problemas de
programación y la métrica \emph{pass@k}, que estima la probabilidad de que al
menos una de $k$ generaciones supere las pruebas asociadas
\cite{chen2021codex}.

Esa capacidad se ha caracterizado también en el dominio de los sistemas
empotrados. Un estudio experimental sobre desarrollo y depuración concluye que
los modelos actuales generan controladores a nivel de registro y código de
comunicación funcional, y que incluso cuando el resultado no es correcto el
razonamiento que lo acompaña resulta de utilidad para el desarrollador
\cite{englhardt2023embedded}. El mismo trabajo aporta dos observaciones
determinantes para esta pasantía: parte del conocimiento del modelo sobre el
hardware procede de la presencia de hojas de datos en su corpus de
entrenamiento, y suministrarle una representación legible de los registros
pertinentes mejora sustancialmente el resultado.

En paralelo, la técnica de generación aumentada por recuperación, conocida por
sus siglas en inglés RAG (\emph{Retrieval-Augmented Generation}), ha demostrado
que combinar un modelo de lenguaje con un sistema de recuperación sobre un
corpus externo permite anclar las respuestas a documentos concretos y
actualizarlas sin reentrenar el modelo \cite{lewis2020rag}. Su aplicación a la
generación de firmware ya ha sido explorada: un trabajo reciente la emplea para
producir capas de abstracción de hardware para microcontroladores de otro
fabricante, integrándola en el flujo de desarrollo existente
\cite{haug2025microcontrollers}.

Existe, por tanto, evidencia de que la técnica es viable y de que el problema
concreto —la fidelidad a los registros de un dispositivo— es precisamente el que
la recuperación puede resolver. Lo que no consta es su realización sobre la
familia PIC18-Q10 de Microchip y su cadena de herramientas, ni un ciclo de
validación que verifique automáticamente el código producido.

\subsection*{Justificación e importancia}

La utilidad de un asistente de generación de código para firmware depende por
completo de la exactitud de lo que produce. Un nombre de registro inexistente,
un campo de bits con anchura equivocada o la omisión de un paso de
inicialización no generan una respuesta aproximada, sino un programa que no
compila o, peor, uno que compila y no funciona. En este dominio no existe la
respuesta parcialmente correcta.

Los modelos de lenguaje presentan tres limitaciones que agravan esa exigencia.
La primera es la generación de contenido incorrecto con apariencia plausible. La
segunda es el corte de conocimiento: el modelo solo dispone de lo que estaba en
su corpus de entrenamiento, sin acceso a documentación posterior ni a las
erratas que el fabricante publica para corregir sus propios manuales. La
tercera, específica de este ámbito, es la escasez de vocabulario especializado:
los nombres de registro y las constantes de hardware aparecen con muy baja
frecuencia en corpus generalistas, de modo que el modelo dispone de poca
evidencia estadística sobre ellos.

La recuperación aumentada aborda las tres simultáneamente. El conocimiento del
dispositivo deja de residir en los parámetros del modelo y pasa a un índice
consultable: se actualiza reindexando, la respuesta queda anclada a un fragmento
verificable y los identificadores exactos llegan al modelo en el propio
contexto. A ello se añade una ventaja que resulta decisiva en documentación
técnica: la jerarquía entre fuentes puede imponerse de forma explícita al
recuperar, de manera que la errata de silicio prevalezca sobre el manual al que
corrige \cite{microchip2024errata}. Un modelo ajustado sobre ambos documentos
los mezclaría en sus pesos sin conservar esa relación de precedencia.

Para Domotai C.A., la relevancia es directa: reducir el tiempo que media entre
la elección de un microcontrolador y la obtención de firmware funcional,
conservando la trazabilidad de cada decisión de configuración hasta la sección
del manual que la respalda.

\subsection*{Planteamiento del problema}

Configurar un periférico en la familia PIC18-Q10 no se agota en escribir los
registros del propio módulo. Intervienen al menos tres mecanismos transversales
que actúan sobre él y que se documentan en secciones alejadas del manual: la
selección de pines de periférico, que desacopla las señales digitales de los
pines físicos y exige una secuencia de desbloqueo; la deshabilitación selectiva
de módulos, que hace que un periférico no responda a ninguna escritura; y el bit
que habilita el juego de instrucciones extendido, que altera la semántica del
direccionamiento de datos para todo el programa. A ellos se suma el
direccionamiento por bancos de la memoria de datos.

La consecuencia es que un fragmento de código correcto según la sección del
periférico puede fallar por causas descritas doscientas páginas más adelante. La
información necesaria para escribir una inicialización completa está
distribuida, no concentrada.

Esa distribución convive con un segundo obstáculo. La hoja de datos del
dispositivo supera las setecientas páginas \cite{microchip2024ds40002043}, de
modo que no cabe íntegra en la ventana de contexto de los modelos de uso
corriente; y aun cuando cupiera, la evidencia experimental muestra que el
aprovechamiento del contexto se degrada con su longitud, y que los elementos
situados en posiciones centrales se utilizan peor que los de los extremos
\cite{liu2024lost}. Suministrar el manual completo no es, por tanto, una
alternativa viable a recuperar selectivamente lo pertinente.

A ello se añade que la información determinante —mapas de registros, asignación
de pines, mapas de memoria— reside en tablas y figuras, no en prosa continua, lo
que condiciona tanto la extracción del documento como su segmentación.

El problema que este trabajo aborda es, en consecuencia, el de construir un
sistema capaz de recuperar de esa documentación los fragmentos pertinentes a una
petición concreta, entregárselos a un modelo de lenguaje y verificar
automáticamente que el código resultante es correcto para el dispositivo.

\subsection*{Objetivo general}

Desarrollar un módulo de generación de código asistido por un modelo de lenguaje
de gran escala que, mediante un sistema de generación aumentada por
recuperación, produzca código verificable para el microcontrolador de 8 bits
PIC18F47Q10 de Microchip, integrando su hoja de datos, la documentación de los
periféricos y un conjunto de configuraciones base validadas.

\subsection*{Objetivos específicos}

\begin{enumerate}
  \item \textbf{Revisar el estado del arte.} Analizar las técnicas actuales de
        generación de código asistida por inteligencia artificial y de
        generación aumentada por recuperación, identificando sus capacidades,
        sus limitaciones y las necesidades del problema a resolver.

  \item \textbf{Configurar el entorno de desarrollo.} Establecer la cadena de
        herramientas para el PIC18F47Q10 —el entorno de desarrollo de Microchip,
        el compilador XC8 y el configurador de código MCC—, adquirir el manejo
        de la tarjeta de evaluación correspondiente y definir los criterios de
        eficiencia en el uso de la memoria de programa y de datos.

  \item \textbf{Generar las configuraciones y ejemplos base.} Desarrollar,
        compilar y validar en hardware un conjunto de ejemplos de referencia
        para los periféricos del dispositivo: pines de entrada y salida de
        propósito general, temporizadores, modulación por ancho de pulso,
        convertidor analógico-digital, convertidor digital-analógico,
        comparadores y las interfaces de comunicación EUSART,
        I\textsuperscript{2}C y SPI. El proceso se acompañará de la
        documentación de los registros y las restricciones de cada periférico.

  \item \textbf{Construir la base de conocimiento y el sistema de
        recuperación.} Procesar y segmentar la hoja de datos y la documentación
        de los periféricos, e indexar los ejemplos validados, de forma que se
        permita la recuperación precisa de la información técnica pertinente.

  \item \textbf{Desarrollar el módulo de generación.} Integrar el modelo de
        lenguaje con el sistema de recuperación, diseñando las instrucciones de
        sistema y el flujo que conduzca a una generación de código efectiva.

  \item \textbf{Evaluar la validación automática del código.} Establecer un
        mecanismo de compilación, verificación y realimentación del código
        generado, incorporando un ciclo que reinyecte los errores del compilador
        al modelo para su corrección iterativa.

  \item \textbf{Evaluar el desempeño y elaborar la documentación técnica.}
        Definir y medir métricas de calidad —tasa de compilación exitosa,
        fidelidad a las configuraciones base y eficiencia en recursos—,
        comparando el desempeño del sistema con y sin recuperación, y generar la
        documentación con la arquitectura, los resultados y las especificaciones
        del módulo.
\end{enumerate}

\subsection*{Estructura del informe}

El documento se organiza en cuatro capítulos. El primero describe la empresa
donde se desarrolló la pasantía. El segundo reúne el marco teórico, recorriendo
desde la arquitectura del microcontrolador y su cadena de herramientas hasta los
modelos de lenguaje, la generación aumentada por recuperación y los criterios de
evaluación. El tercero expone la metodología seguida, y el cuarto presenta los
resultados obtenidos junto con su discusión. Cierran el documento las
conclusiones y recomendaciones, las referencias consultadas y los apéndices.
